\section{Results}

\subsection{RQ1: Cheap Evidence Preserves Energy Decisions}

Table~\ref{tab:evidence-results} summarizes the two sealed holdouts.  On the
eight-workload v2 holdout, energy MAPE is \CalibratedEnergyMAPE{}, median APE is
\EnergyMedianAPE{}, and the maximum APE is \EnergyMaxAPE{}.  Energy rank
correlation is \EnergySpearman{}, and \EnergyPairwise{} pairwise orderings
(\EnergyPairwisePct{}) are correct.  Median repeat CV is only
\EnergyRepeatCV{}, indicating that prediction error, rather than measurement
noise, dominates.  Throughput and TPOT MAPE are
\CalibratedThroughputMAPE{} and \CalibratedTPOTMAPE{}.  TTFT is less accurate
at \CalibratedTTFTMAPE{}, motivating the separate release gate below.

Without refitting, scope confirmation obtains \ConfirmationEnergyMAPE{}
energy MAPE, \ConfirmationEnergySpearman{} rank correlation, and
\ConfirmationEnergyPairwise{} correct pairwise orderings.  Its largest energy
error is \ConfirmationEnergyMaxAPE{}, so the preregistered 15\% primary
endpoint passes on all nine workloads.  The result supports CPU-first energy
screening inside the measured workload and runtime envelope; it does not
establish cross-GPU or multi-node accuracy.

\begin{table}[t]
\centering
\caption{Frozen-model results on two disjoint H100 holdouts.  Each workload
has three repeats; MAPE compares the frozen prediction with the observed
median.}
\label{tab:evidence-results}
\small
\setlength{\tabcolsep}{3.8pt}
\begin{tabular}{lrrrr}
\toprule
Study & Workloads & Runs & Energy MAPE & $\rho$ \\
\midrule
Holdout v2 & \HoldoutWorkloads{} & \HoldoutMeasurements{} &
\CalibratedEnergyMAPE{} & \EnergySpearman{} \\
Scope v3 & \ConfirmationWorkloads{} & \ConfirmationMeasurements{} &
\ConfirmationEnergyMAPE{} & \ConfirmationEnergySpearman{} \\
\bottomrule
\end{tabular}
\end{table}


\subsection{RQ2: Scope Gating Prevents a False Latency Claim}

Aggregate TTFT MAPE in the confirmation is \ConfirmationTTFTMAPE{}, which by
itself would reject a global TTFT claim.  Stratification reveals the boundary
in Table~\ref{tab:scope-results}.  At concurrency four, TTFT MAPE is
\KFourTTFTMAPE{} and maximum APE is \SupportedTTFTMaxAPE{}, passing the
preregistered 20\% rule.  At concurrency one and two, TTFT MAPE rises to
\KOneTTFTMAPE{} and \KTwoTTFTMAPE{}; the combined sparse-stratum error is
\SparseTTFTMAPE{}.  The generated decision therefore supports latency at
concurrency at least four and abstains below four.  Energy remains stable in
all three strata.  This is the desired behavior for an evidence-gated agent:
it narrows a claim instead of turning confident energy estimates into an
unsupported latency recommendation.

\begin{table}[t]
\centering
\caption{Preregistered scope-confirmation MAPE by concurrency.}
\label{tab:scope-results}
\small
\setlength{\tabcolsep}{4.2pt}
\begin{tabular}{crrrr}
\toprule
Conc. & Energy & Throughput & TPOT & TTFT \\
\midrule
1 & \KOneEnergyMAPE{} & \KOneThroughputMAPE{} &
\KOneTPOTMAPE{} & \KOneTTFTMAPE{} \\
2 & \KTwoEnergyMAPE{} & \KTwoThroughputMAPE{} &
\KTwoTPOTMAPE{} & \KTwoTTFTMAPE{} \\
4 & \KFourEnergyMAPE{} & \KFourThroughputMAPE{} &
\KFourTPOTMAPE{} & \KFourTTFTMAPE{} \\
\bottomrule
\end{tabular}
\end{table}


\subsection{RQ3: Configuration-Search Efficiency}

The runtime can already execute sealed replay episodes for IPIG, random,
cost-blind, and cheapest-first policies and produce all metrics in
Table~\ref{tab:agent-results}.  The bundled three-candidate corpus is explicitly
synthetic and is used only for software tests.  Publication values require the
frozen H100 serving-configuration corpus described in Section~VII; until that
campaign is sealed, we make no claim that IPIG outperforms a search baseline.

\begin{table}[t]
\centering
\caption{Sealed controller results.  Policy rows use
\BudgetEpisodesPerPoint{} matched episodes/budget; AUC spans \BudgetLevels{}
budgets.  Planner rows use \PlannerCallsPerSet{} calls/set; V2 is a disjoint
holdout.  Guard event is V1 fallback or V2 intervention; latency is P50/P95 ms.}
\label{tab:agent-results}
\footnotesize
\setlength{\tabcolsep}{3.0pt}
\textbf{(a) Acquisition policy}\par\vspace{2pt}
\begin{tabular}{lrrr}
\toprule
Policy & Hit@.12 & Mean GPU-h & Success AUC \\
\midrule
\ipig{} & \IPIGOracleHit{} & \IPIGGPUHours{} & \IPIGSuccessAUC{} \\
Random & \RandomOracleHit{} & \RandomGPUHours{} & \RandomSuccessAUC{} \\
Cost-blind & \CostBlindOracleHit{} & \CostBlindGPUHours{} & \CostBlindSuccessAUC{} \\
Cheapest-first & \CheapestOracleHit{} & \CheapestGPUHours{} & \CheapestSuccessAUC{} \\
\bottomrule
\end{tabular}
\par\vspace{2pt}
\textbf{(b) Semantic planner}\par\vspace{2pt}
\begin{tabular}{lrrrr}
\toprule
Set & Raw & Admitted & Guard event & P50/P95 \\
\midrule
V1 development & \PlannerVoneRaw{} & \PlannerVoneAdmitted{} &
\PlannerVoneFallback{} & \PlannerVoneMedianLatency{}/\PlannerVonePninetyfiveLatency{} \\
V2 holdout & \PlannerVtwoRaw{} & \PlannerVtwoAdmitted{} &
\PlannerVtwoIntervention{} & \PlannerVtwoMedianLatency{}/\PlannerVtwoPninetyfiveLatency{} \\
\bottomrule
\end{tabular}
\end{table}


\subsection{RQ4: Bounded Runtime Behavior}

All 58 automated tests pass.  Injected executor failure becomes a failed
evidence record charged at its declared cost, and the next event enters
\textsc{Repair}.  Malformed LLM output triggers a recorded rule-based fallback.
Replay and topology tools cannot relabel L0/L2 outputs as verified, budget
overruns raise a hard error, and the returned recommendation set is computed
only from successful L4 evidence that satisfies the SLO.  These tests establish
mechanism enforcement, not the semantic quality of a particular planner model.
